\scanpage{36}
\mspage{\handwritten{4}}
% archive sheet number: 35
\uline{présentation finie, et $R_j$ une relation d'équivalence sur $X_j$ définie par un $S_j$-monomorphisme $v_j:R_j\handwritten{\longrightarrow}X_j\handwritten{\times}_{S_j}X_j$ tel que le morphisme $\operatorname{pr}_{1,j}\mathbin{o} v_j:R_j\handwritten{\longrightarrow}X_j$ déduit de $v_j$ par composition avec la première projection $X_j\handwritten{\times}_{S_j}X_j\handwritten{\unsure{\longrightarrow|\longleftrightarrow}}X_j$ soit plat et de présentation finie. Alors, pour que le faisceau quotient $X_j\handwritten{\times}_{S_j}S/R_j\handwritten{\times}_{S_j}S$ pour la topologie $\handwritten{\mathbb{T}}$ soit représentable, il faut et il suffit qu'il existe $i\handwritten{\in}I$, $i\handwritten{\geqslant}j$, le faisceau quotient $X_j\handwritten{\times}_{S_j}S_i/R_j{}_{S_j}S_i$ pour la topologie $\handwritten{\mathbb{T}}$ soit représentable.}

Compte tenu des énoncés de (EGA IV 8.8.2, 8.8.3, 8.10.5 et 11.2.6) rappelés en (10.0), ce lemme est conséquence\deleted{\unsure{du|un}} du résultat suivant :

\uline{Lemme 10.8.}- \uline{Soit $\handwritten{\mathbb{T}}$ la topologie fppf ou fpqc ; soient $\deleted{\unsure{j\in I,}}$ $X$ un $S$-préschéma de présentation finie (resp. localement de présentation finie), $R$ une relation d'\handwritten{é}quivalence sur $X$ définie par un monomorphisme $v:R\handwritten{\longrightarrow}X\handwritten{\times}_S X$ tel que $\operatorname{pr}_1\mathbin{o} v$ soit plat et de présentation finie (resp. plat et localement de présentation finie). Alors les conditions suivantes sont équivalentes :}

\uline{(i) Le faisceau quotient $X/R$ pour la topologie $\handwritten{\mathbb{T}}$ est représentable.}

\uline{(ii) il existe un $S$-préschéma de présentation finie (resp. localement de présentation finie ) $Y$ et un morphisme fidèlement plat $p:X\handwritten{\longrightarrow}Y$ tel que le diagramme}

\diagram{\(\text{(D)}\)\qquad\begin{tikzcd}[ampersand replacement=\&,column sep=large,row sep=large]
R \handwritten{\arrow[r,"\operatorname{pr}_1\mathbin{o} v"]} \handwritten{\arrow[d,"\operatorname{pr}_2\mathbin{o} v"']} \& X \handwritten{\arrow[d,"p"]} \\
X \handwritten{\arrow[r,"p"']} \& X
\end{tikzcd}}

(où $\operatorname{pr}_1$ et $\operatorname{pr}_2$ sont les projections $X\handwritten{\times}_S X\handwritten{\unsure{\substack{\longrightarrow\\[-.6ex]\rightharpoonup}|\rightrightarrows}}X$) \uline{soit cartésien.}

Notons d'abord que d'après (Exp. IV 3.3.2 et 4.4.3), pour que le faisceau $X/R$ \deleted{\illegible{1}} pour la topologie $\handwritten{\mathbb{T}}$ soit représentable par $Y$, il faut et il suffit que le diagramme (D) soit cartésien et que $p$ soit couvrant pour la topologie $\handwritten{\mathbb{T}}$.

Montrons que i) entraîne ii). L'hypothèse (i) implique que le diagramme (D) est cartésien, donc que $\operatorname{pr}_1\mathbin{o} v$ se déduit de $p$ par changement de base par $p$, et que $p$ est
